% ============================================================
% ESPECIFICACIÓN DE REQUERIMIENTOS DE SOFTWARE — IEEE 830
% Sistema Genérico (Proyecto PW)
% Documento maestro: compilar este archivo en Overleaf (pdfLaTeX)
% ============================================================
\documentclass[12pt,a4paper,oneside]{report}

% ── Codificación e idioma ───────────────────────────────────
\usepackage[utf8]{inputenc}
\usepackage[T1]{fontenc}
\usepackage[spanish,es-tabla]{babel}

% ── Página y tipografía ─────────────────────────────────────
\usepackage[margin=2.5cm,headheight=15pt]{geometry}
\usepackage{setspace}
\onehalfspacing
\usepackage[table]{xcolor}
\usepackage{graphicx}
\usepackage{booktabs}
\usepackage{longtable}
\usepackage{tabularx}
\usepackage{array}
\usepackage{enumitem}
\setlist{nosep}
\usepackage{fancyhdr}
\usepackage{titlesec}

% ── Colores corporativos (paleta del sistema) ───────────────
\definecolor{azulcorp}{HTML}{3A7BD5}
\definecolor{moradocorp}{HTML}{5A4FCF}
\definecolor{grisfila}{HTML}{F1F5F9}

% ── Encabezado y pie de página ──────────────────────────────
\pagestyle{fancy}
\fancyhf{}
\fancyhead[L]{\small\itshape ERS — Sistema Genérico}
\fancyhead[R]{\small\itshape IEEE 830}
\fancyfoot[C]{\thepage}
\renewcommand{\headrulewidth}{0.4pt}

% ── Estilo de capítulos ─────────────────────────────────────
\titleformat{\chapter}[hang]
  {\normalfont\huge\bfseries\color{azulcorp}}
  {\thechapter.}{0.6em}{}
\titlespacing*{\chapter}{0pt}{0pt}{24pt}

% ── Hipervínculos y metadatos ───────────────────────────────
\usepackage[hidelinks]{hyperref}
\hypersetup{
  pdftitle={ERS IEEE 830 - Sistema Genérico},
  pdfauthor={David Bonilla},
  pdfsubject={Especificación de Requerimientos de Software},
  pdfkeywords={IEEE 830, RBAC, requerimientos, Sistema Genérico}
}

% ── Datos editables del documento ───────────────────────────
\newcommand{\nombreSistema}{Sistema Genérico}
\newcommand{\autorDoc}{David Bonilla}
\newcommand{\versionDoc}{1.0}
\newcommand{\cursoDoc}{NRC 30713 — Programación Web}

% ── Ficha de requisito / caso de uso (tabla que corta página) ──
% Uso:
% \begin{ficha}{RF-001 — Nombre del requisito}
%   Descripción  & ...texto... \\ \hline
%   Prioridad    & Alta        \\ \hline
% \end{ficha}
\newenvironment{ficha}[1]{%
  \begin{longtable}{|>{\bfseries\raggedright\arraybackslash}p{3.2cm}|p{11.4cm}|}
  \hline
  \multicolumn{2}{|l|}{\cellcolor{azulcorp!15}\textbf{#1}}\\ \hline
  \endfirsthead
  \hline
  \multicolumn{2}{|l|}{\cellcolor{azulcorp!15}\textbf{#1} \textit{(continuación)}}\\ \hline
  \endhead
}{%
  \end{longtable}
}

% ============================================================
\begin{document}

% ── Portada ─────────────────────────────────────────────────
\begin{titlepage}
    \centering
    \vspace*{1.5cm}

    {\color{azulcorp}\rule{\textwidth}{1.5pt}}

    \vspace{1.2cm}
    {\Huge\bfseries Especificación de Requerimientos\\[0.35cm] de Software\par}

    \vspace{0.8cm}
    {\Large Estándar IEEE 830\par}

    \vspace{1.6cm}
    {\color{moradocorp}\Huge\bfseries \nombreSistema\par}
    \vspace{0.4cm}
    {\large Sistema web de administración con control de acceso\\ basado en roles (RBAC)\par}

    \vspace{1.2cm}
    {\color{azulcorp}\rule{0.6\textwidth}{0.8pt}}

    \vfill

    \begin{tabular}{ll}
        \textbf{Autor:}    & \autorDoc \\[0.15cm]
        \textbf{Curso:}    & \cursoDoc \\[0.15cm]
        \textbf{Versión:}  & \versionDoc \\[0.15cm]
        \textbf{Fecha:}    & \today \\
    \end{tabular}

    \vspace{1.5cm}
    {\color{azulcorp}\rule{\textwidth}{1.5pt}}
\end{titlepage}


\tableofcontents
\clearpage

% ============================================================
% CAPÍTULO 1 — INTRODUCCIÓN (IEEE 830 §1)
% ============================================================
\chapter{Introducción}

Este documento constituye la Especificación de Requerimientos de Software (ERS)
del \textbf{\nombreSistema}, elaborada conforme al estándar \textbf{IEEE 830}.
Todos los requerimientos, reglas de negocio y casos de uso aquí descritos
fueron extraídos directamente del código fuente y de la base de datos del
proyecto, por lo que reflejan con exactitud el comportamiento real del sistema
implementado.

\section{Propósito}

El propósito de este documento es especificar de forma completa, precisa y
verificable los requerimientos funcionales y no funcionales del
\nombreSistema, un sistema web de administración con control de acceso basado
en roles (RBAC). El documento está dirigido a:

\begin{itemize}
    \item El equipo docente y evaluador del curso \cursoDoc, como evidencia
          formal del alcance implementado.
    \item Los desarrolladores actuales y futuros del sistema, como referencia
          de mantenimiento y evolución.
    \item Cualquier interesado que requiera comprender el funcionamiento del
          sistema sin examinar su código fuente.
\end{itemize}

\section{Alcance}

El producto descrito se denomina \textbf{\nombreSistema} (nombre interno del
proyecto: \textit{Proyecto\_pw}). Es una aplicación web de tipo página única
(SPA) con arquitectura \textit{shell + frames} que provee:

\begin{itemize}
    \item Autenticación de usuarios con captcha propio, bloqueo por intentos
          fallidos y sesiones con token de expiración.
    \item Un sistema RBAC de dos niveles: acceso a pantallas (\textit{Frames})
          y acciones internas dentro de cada pantalla, configurable por rol
          desde la interfaz.
    \item Administración de usuarios, roles y permisos (exclusiva del
          Administrador).
    \item Un menú de navegación personalizable por cada usuario mediante
          arrastrar y soltar (SuperMenus propios), con disponibilidad global
          de opciones controlada por el Administrador.
    \item Un Dashboard con accesos rápidos personalizables, estadísticas
          (solo Administrador) y frases motivadoras (solo usuarios estándar).
    \item Gestión del perfil propio (consulta de datos y cambio de contraseña).
    \item Cinco Frames de ejemplo que demuestran el sistema de permisos,
          incluido un CRUD de tareas con permisos por operación.
\end{itemize}

Queda fuera del alcance del sistema: la recuperación de contraseñas por
correo, el registro público de usuarios (las cuentas solo las crea el
Administrador), la internacionalización (la interfaz es únicamente en
español) y la exposición de una API pública para terceros.

\section{Definiciones, acrónimos y abreviaturas}

\begin{longtable}{|>{\bfseries}p{3.6cm}|p{11cm}|}
\hline
\rowcolor{azulcorp!15} \textbf{Término} & \textbf{Definición} \\ \hline
\endhead
Frame & Pantalla del sistema que se carga dentro del marco de contenido del
        shell (por ejemplo: Usuarios, Dashboard, Frame 1). \\ \hline
ItemMenu & Opción de menú que abre un Frame. Es global al sistema y se
        registra en la tabla \texttt{menu}. \\ \hline
SuperMenu & Contenedor de ItemMenus creado por un usuario para organizar su
        menú personal. Es individual: cada usuario posee los suyos. \\ \hline
Shell & Página contenedora única (\texttt{dashboard/index.html}) que aloja la
        barra lateral, la barra superior y el marco donde se cargan los
        Frames sin recargar la estructura. \\ \hline
Módulo & Identificador RBAC asociado a un Frame (p.~ej. \texttt{usuarios},
        \texttt{frame1}), sobre el que se otorgan permisos. \\ \hline
Acción & Operación autorizable dentro de un módulo: \texttt{leer},
        \texttt{crear}, \texttt{editar} o \texttt{eliminar}. \\ \hline
Registro de Frames & Catálogo central del servidor
        (\texttt{servidor/permisos/registro.php}) que define los módulos
        existentes y sus acciones internas. \\ \hline
Semilla de menú & Estructura de menú inicial que se genera automáticamente
        una única vez al crear un usuario. \\ \hline
Token de sesión & Cadena hexadecimal de 64 caracteres generada
        aleatoriamente que identifica una sesión activa. \\ \hline
Captcha & Verificación gráfica generada en un lienzo (canvas) del navegador
        con seis caracteres aleatorios. \\ \hline
RBAC & \textit{Role-Based Access Control}: control de acceso basado en roles. \\ \hline
ERS & Especificación de Requerimientos de Software. \\ \hline
SPA & \textit{Single Page Application}: aplicación de página única. \\ \hline
CRUD & Crear, Leer, Actualizar y Eliminar (\textit{Create, Read, Update,
        Delete}). \\ \hline
PDO & \textit{PHP Data Objects}: capa de acceso a datos de PHP. \\ \hline
JSON & \textit{JavaScript Object Notation}: formato de intercambio de datos. \\ \hline
DnD & \textit{Drag \& Drop}: interacción de arrastrar y soltar. \\ \hline
bcrypt & Algoritmo de hash utilizado para almacenar contraseñas. \\ \hline
\end{longtable}

\section{Referencias}

\begin{itemize}
    \item IEEE Std 830-1998, \textit{IEEE Recommended Practice for Software
          Requirements Specifications}.
    \item Código fuente del proyecto \textit{Proyecto\_pw} (carpetas
          \texttt{paginas/}, \texttt{scripts/}, \texttt{servidor/},
          \texttt{estilos/}).
    \item Volcado de la base de datos \texttt{proyecto\_pw} (MySQL 8.0),
          archivo \texttt{proyecto\_pw.sql}.
    \item Documentación oficial de PHP 7.3 y MySQL 8.0.
\end{itemize}

\section{Visión general del documento}

El resto del documento se organiza de la siguiente manera: el
Capítulo~\ref{cap:descripcion} presenta la descripción general del producto,
su arquitectura, funciones, usuarios, restricciones y dependencias; el
Capítulo~\ref{cap:rf} especifica los requerimientos funcionales agrupados por
módulo; el Capítulo~\ref{cap:rnf} recoge los requerimientos no funcionales;
el Capítulo~\ref{cap:cu} desarrolla los casos de uso; el
Capítulo~\ref{cap:rn} enumera las reglas de negocio implementadas con su
evidencia en el código; finalmente, el Capítulo~\ref{cap:anexos} contiene los
anexos técnicos (modelo de datos, matriz de permisos, catálogo de endpoints y
estructura del proyecto).

% ============================================================
% CAPÍTULO 2 — DESCRIPCIÓN GENERAL (IEEE 830 §2)
% ============================================================
\chapter{Descripción general}\label{cap:descripcion}

\section{Perspectiva del producto}

El \nombreSistema{} es un producto autónomo (no forma parte de un sistema
mayor) construido sobre una arquitectura web cliente--servidor de tres capas:

\begin{description}
    \item[Capa de presentación (cliente).] HTML5, CSS3 y JavaScript sin
    frameworks. La navegación sigue el patrón \textit{shell + frames}: existe
    una única página contenedora (\texttt{paginas/dashboard/index.html}) que
    renderiza una sola vez la barra lateral de navegación y la barra
    superior, y carga cada Frame dentro de un marco interno
    (\texttt{iframe}) sin recargar la estructura general. El estilo visual
    (\textit{glassmorphism}: paneles flotantes translúcidos con desenfoque
    sobre una imagen de fondo con capa oscura) se hereda de la pantalla de
    inicio de sesión hacia todo el sistema.

    \item[Capa de lógica (servidor).] PHP~7.3 con acceso a datos mediante
    PDO. Cada operación de negocio se expone como un endpoint independiente
    que recibe y responde JSON, siguiendo el patrón uniforme: verificar
    sesión $\rightarrow$ verificar permiso RBAC $\rightarrow$ verificar
    disponibilidad global del módulo $\rightarrow$ ejecutar consulta
    $\rightarrow$ responder \texttt{\{ok, msg, data\}}. Las bibliotecas
    compartidas son \texttt{servidor/config.php} (conexión, sesión,
    autorización), \texttt{servidor/permisos/registro.php} (catálogo central
    de módulos y acciones) y \texttt{servidor/menu/semilla.php} (menú
    inicial de usuarios nuevos).

    \item[Capa de datos.] Base de datos MySQL~8.0 (\texttt{proyecto\_pw})
    con 11 tablas, integridad referencial completa mediante claves foráneas
    y restricciones de unicidad sobre las claves de negocio.
\end{description}

El estado de sesión del cliente (token, usuario, permisos, menú) se conserva
en \texttt{sessionStorage} y se revalida contra el servidor en cada carga de
página.

\subsection{Componentes del cliente}

\begin{longtable}{|>{\ttfamily}p{4.2cm}|p{10.4cm}|}
\hline
\rowcolor{azulcorp!15} \textbf{Archivo} & \textbf{Responsabilidad} \\ \hline
\endhead
config.js & Constantes globales, rutas de páginas y catálogo de endpoints. \\ \hline
utils.js & Peticiones JSON, alertas flotantes, indicador de carga, formato y
           saneamiento de texto (prevención XSS). \\ \hline
session.js & Lectura/escritura del estado de sesión y verificación de
             permisos en el cliente. \\ \hline
router.js & Guardias de navegación: protección de páginas, verificación de
            permisos y de disponibilidad del módulo, redirecciones. \\ \hline
validaciones.js & Validación de formularios y política de contraseñas. \\ \hline
shell.js & Controlador del shell: carga de Frames en el marco, título
           dinámico, ítem activo, enlace profundo por \textit{hash}. \\ \hline
menu.js & Renderizado de la barra lateral (acordeones) y barra superior. \\ \hline
auth.js & Inicio de sesión, captcha, bloqueo por intentos, mostrar/ocultar
          contraseña, cierre de sesión. \\ \hline
dashboard.js, usuarios.js, roles.js, permisos.js, menu-personal.js,
menu-config.js, cambiar-password.js, frame1.js & Controlador de cada Frame. \\ \hline
\end{longtable}

\section{Funciones del producto}

El sistema ofrece doce pantallas (Frames) organizadas en los siguientes
grupos funcionales:

\begin{longtable}{|p{3.4cm}|p{11.2cm}|}
\hline
\rowcolor{azulcorp!15} \textbf{Grupo} & \textbf{Funciones} \\ \hline
\endhead
Autenticación & Inicio de sesión con captcha y bloqueo por intentos; cierre
                de sesión; expiración y verificación de sesión; cambio de
                contraseña obligatorio en el primer acceso. \\ \hline
Navegación & Shell con menú lateral filtrado por permisos, disponibilidad
             global y organización personal; carga de Frames sin recarga;
             modo móvil. \\ \hline
Dashboard & Saludo contextual, accesos rápidos personalizables por usuario,
            estadísticas del sistema (solo Administrador) y frase motivadora
            aleatoria (solo usuarios estándar). \\ \hline
Administración & Gestión de usuarios (listar, filtrar, crear, editar,
                 activar/desactivar), de roles (listar, crear, editar) y de
                 permisos por rol en dos niveles (acceso a Frames y acciones
                 internas). Exclusiva del Administrador. \\ \hline
Menús & Personalización individual del menú mediante arrastrar y soltar
        (SuperMenus propios) y administración de la disponibilidad global de
        ItemMenus (solo Administrador). \\ \hline
Mi Perfil & Consulta de los datos propios y cambio de contraseña con
            indicador de fortaleza. \\ \hline
Frames de ejemplo & Frame~1: CRUD de tareas con permisos por operación;
                    Frames~2--5: pantallas de demostración con control de
                    acceso. \\ \hline
\end{longtable}

\section{Características de los usuarios}

\begin{longtable}{|p{2.6cm}|p{5.6cm}|p{5.9cm}|}
\hline
\rowcolor{azulcorp!15} \textbf{Actor} & \textbf{Descripción} &
\textbf{Privilegios} \\ \hline
\endhead
Visitante & Persona no autenticada. No requiere conocimientos previos. &
  Únicamente puede acceder a la pantalla de inicio de sesión; cualquier otra
  página lo redirige al login. \\ \hline
Usuario estándar & Usuario autenticado con alguno de los roles Rol1--Rol7.
  Usuario final con conocimientos básicos de navegación web. &
  Accede solo a los Frames habilitados a su rol; siempre dispone de
  Dashboard, Mi Perfil, la lista de usuarios en modo lectura y la
  personalización de su propio menú; dentro de cada Frame solo ve las
  acciones autorizadas. \\ \hline
Administrador & Usuario con el rol \texttt{Admin} (identificador de rol~1).
  Responsable de la operación del sistema. &
  Posee implícitamente todos los permisos (no es configurable); es el único
  que administra usuarios, roles, permisos y la disponibilidad global de los
  menús. \\ \hline
\end{longtable}

\section{Restricciones}

\begin{itemize}
    \item El servidor requiere PHP~7.3 o superior con PDO/MySQL y una base
          de datos MySQL~8.0 o MariaDB equivalente.
    \item El cliente no utiliza frameworks ni bibliotecas externas
          (JavaScript puro, CSS propio); no existen dependencias de
          terceros ni gestores de paquetes.
    \item Toda la comunicación cliente--servidor se realiza por HTTP con
          peticiones POST y cuerpo JSON; el token de sesión viaja en el
          cuerpo de cada petición.
    \item Las contraseñas se almacenan exclusivamente con hash bcrypt
          (costo 12); ninguna operación devuelve la contraseña.
    \item La autorización se valida siempre en el servidor; las
          comprobaciones del cliente son únicamente de experiencia de
          usuario.
    \item Los módulos del sistema deben estar declarados en el registro
          central de Frames; el guardado de permisos rechaza módulos o
          acciones no registrados.
    \item Las rutas del proyecto están ancladas al prefijo
          \texttt{/NRC30713-Web/Proyecto\_pw/} del servidor web.
\end{itemize}

\section{Suposiciones y dependencias}

\begin{itemize}
    \item Se asume un servidor Apache local (entorno de referencia: AppServ
          en el puerto 8080) con la base de datos \texttt{proyecto\_pw}
          creada e inicializada con los datos semilla.
    \item Se asume el uso de un navegador moderno con JavaScript,
          \texttt{sessionStorage} e \texttt{iframe} habilitados.
    \item La personalización del menú mediante arrastrar y soltar utiliza la
          API nativa de HTML5 \textit{Drag \& Drop}, orientada a dispositivos
          con puntero (escritorio).
    \item La existencia del usuario administrador inicial
          (\texttt{admin}, rol~1) es una precondición de operación: el
          sistema no ofrece registro público de cuentas.
    \item Los relojes del servidor determinan la expiración de sesiones y
          bloqueos; se asume una zona horaria consistente.
\end{itemize}

% ============================================================
% CAPÍTULO 3 — REQUERIMIENTOS FUNCIONALES (IEEE 830 §3)
% 54 requerimientos extraídos del código fuente, agrupados por módulo.
% ============================================================
\chapter{Requerimientos funcionales}\label{cap:rf}

Los requerimientos se agrupan por módulo funcional. La prioridad se clasifica
en \textbf{Alta} (esencial para operar el sistema), \textbf{Media}
(funcionalidad relevante no crítica) y \textbf{Baja} (mejora de experiencia).

% ────────────────────────────────────────────────────────────
\section{Autenticación y sesión}

\begin{ficha}{RF-001 — Inicio de sesión}
Descripción & El sistema permite autenticarse con nombre de usuario y
contraseña, verificadas contra el hash bcrypt almacenado. \\ \hline
Prioridad & Alta \\ \hline
Actor & Visitante \\ \hline
Entradas & Usuario, contraseña y texto del captcha. \\ \hline
Proceso & Se validan los campos en el cliente; se verifica el captcha; el
servidor comprueba el bloqueo por intentos, busca al usuario, verifica la
contraseña con bcrypt, valida el estado del usuario y de su rol, genera el
token de sesión y devuelve usuario, permisos y token. \\ \hline
Salidas & Sesión iniciada y redirección al shell; o mensaje de error con los
intentos restantes. \\ \hline
Precondiciones & Cuenta existente y activa; captcha vigente; sin bloqueo
activo. \\ \hline
Postcondiciones & Token registrado en la tabla \texttt{sesiones}; estado de
sesión guardado en el navegador; intentos fallidos reiniciados. \\ \hline
\end{ficha}

\begin{ficha}{RF-002 — Verificación captcha}
Descripción & El login exige transcribir un captcha gráfico de seis
caracteres generado localmente en un lienzo con ruido visual. \\ \hline
Prioridad & Alta \\ \hline
Actor & Visitante \\ \hline
Entradas & Texto transcrito por el usuario. \\ \hline
Proceso & El captcha se genera al cargar la página y puede regenerarse
pulsando la imagen o el botón de refresco; expira a los cinco minutos; la
comparación no distingue mayúsculas de minúsculas. \\ \hline
Salidas & Validación superada, o mensaje de error con regeneración
automática del captcha. \\ \hline
Precondiciones & Página de login cargada. \\ \hline
Postcondiciones & Si es incorrecto o expiró, se genera un captcha nuevo y se
limpia el campo. \\ \hline
\end{ficha}

\begin{ficha}{RF-003 — Bloqueo por intentos fallidos}
Descripción & Tras cinco intentos fallidos consecutivos, el acceso de esa
cuenta se bloquea durante quince minutos. \\ \hline
Prioridad & Alta \\ \hline
Actor & Visitante \\ \hline
Entradas & Intentos de inicio de sesión fallidos. \\ \hline
Proceso & El servidor contabiliza los intentos por usuario e IP
(\texttt{login\_intentos}); el cliente replica el contador y muestra un
aviso con los minutos restantes, actualizado cada segundo, deshabilitando el
botón de ingreso. Un inicio exitoso limpia el contador. \\ \hline
Salidas & Mensaje de bloqueo con tiempo restante; intentos restantes en cada
fallo. \\ \hline
Precondiciones & --- \\ \hline
Postcondiciones & Cuenta bloqueada temporalmente; desbloqueo automático al
cumplirse el tiempo. \\ \hline
\end{ficha}

\begin{ficha}{RF-004 — Mostrar/ocultar contraseña}
Descripción & El campo de contraseña del login incluye un ícono de ojo que
alterna entre texto oculto y visible. \\ \hline
Prioridad & Baja \\ \hline
Actor & Visitante \\ \hline
Entradas & Pulsación sobre el ícono. \\ \hline
Proceso & Se alterna el tipo del campo entre \texttt{password} y
\texttt{text} sin borrar el contenido ni robar el foco del campo. \\ \hline
Salidas & Contraseña visible u oculta; el ícono refleja el estado. \\ \hline
Precondiciones & Página de login cargada. \\ \hline
Postcondiciones & El proceso de autenticación no se ve afectado. \\ \hline
\end{ficha}

\begin{ficha}{RF-005 — Sesión con token de expiración}
Descripción & Cada sesión se identifica con un token aleatorio de 64
caracteres hexadecimales con vigencia de una hora. \\ \hline
Prioridad & Alta \\ \hline
Actor & Sistema \\ \hline
Entradas & Inicio de sesión exitoso. \\ \hline
Proceso & El servidor genera el token con \texttt{random\_bytes}, elimina
las sesiones anteriores del usuario (sesión única) y purga las sesiones
caducadas de todos los usuarios. \\ \hline
Salidas & Token entregado al cliente y almacenado en
\texttt{sessionStorage}. \\ \hline
Precondiciones & Credenciales válidas. \\ \hline
Postcondiciones & Una única sesión vigente por usuario. \\ \hline
\end{ficha}

\begin{ficha}{RF-006 — Protección de páginas y verificación de sesión}
Descripción & Toda página distinta del login verifica la sesión contra el
servidor antes de mostrarse. \\ \hline
Prioridad & Alta \\ \hline
Actor & Sistema \\ \hline
Entradas & Token almacenado en el navegador. \\ \hline
Proceso & El guardia de navegación consulta al servidor; si el token es
inválido, expiró o el usuario fue desactivado, limpia el estado local y
redirige la ventana completa al login; en cada verificación se refrescan los
permisos del rol. \\ \hline
Salidas & Acceso concedido, o redirección al login. \\ \hline
Precondiciones & --- \\ \hline
Postcondiciones & Solo usuarios con sesión válida visualizan contenido. \\ \hline
\end{ficha}

\begin{ficha}{RF-007 — Redirección de usuario autenticado}
Descripción & Si un usuario con sesión vigente abre el login, es redirigido
automáticamente al shell. \\ \hline
Prioridad & Media \\ \hline
Actor & Usuario estándar, Administrador \\ \hline
Entradas & Apertura de la página de login. \\ \hline
Proceso & Se verifica el token contra el servidor; si es válido se reemplaza
la ubicación por el shell del sistema. \\ \hline
Salidas & Navegación directa al Dashboard. \\ \hline
Precondiciones & Sesión vigente. \\ \hline
Postcondiciones & --- \\ \hline
\end{ficha}

\begin{ficha}{RF-008 — Cambio de contraseña obligatorio en el primer acceso}
Descripción & Un usuario con la marca de primer inicio debe cambiar su
contraseña antes de usar el sistema. \\ \hline
Prioridad & Alta \\ \hline
Actor & Usuario estándar, Administrador \\ \hline
Entradas & Primer inicio de sesión de la cuenta. \\ \hline
Proceso & Tras autenticarse, el guardia detecta \texttt{primer\_login} y
redirige al formulario de cambio de contraseña; el botón Cancelar queda
inhabilitado mediante un aviso mientras la marca esté activa. \\ \hline
Salidas & Formulario de cambio de contraseña con aviso de obligatoriedad. \\ \hline
Precondiciones & Cuenta creada por el Administrador o con contraseña
restablecida. \\ \hline
Postcondiciones & Al cambiar la contraseña, la marca se desactiva y el
usuario accede normalmente. \\ \hline
\end{ficha}

\begin{ficha}{RF-009 — Cierre de sesión}
Descripción & El usuario puede cerrar su sesión desde el botón Salir de la
barra superior. \\ \hline
Prioridad & Alta \\ \hline
Actor & Usuario estándar, Administrador \\ \hline
Entradas & Pulsación del botón Salir. \\ \hline
Proceso & Se limpia el estado local, se elimina el token en el servidor y se
redirige al login. \\ \hline
Salidas & Sesión finalizada. \\ \hline
Precondiciones & Sesión activa. \\ \hline
Postcondiciones & El token deja de ser válido. \\ \hline
\end{ficha}

\begin{ficha}{RF-010 — Rechazo de cuentas o roles inactivos}
Descripción & El sistema impide iniciar sesión a usuarios desactivados y a
usuarios cuyo rol esté inactivo. \\ \hline
Prioridad & Alta \\ \hline
Actor & Visitante \\ \hline
Entradas & Credenciales de una cuenta o rol inactivo. \\ \hline
Proceso & Tras validar la contraseña, el servidor consulta el estado del
usuario y el estado de su rol; si alguno es inactivo, se rechaza el acceso
con un mensaje específico. \\ \hline
Salidas & Mensaje ``Usuario inactivo'' o ``Su rol está inactivo''. \\ \hline
Precondiciones & Credenciales correctas. \\ \hline
Postcondiciones & No se crea sesión. \\ \hline
\end{ficha}

% ────────────────────────────────────────────────────────────
\section{Navegación (shell)}

\begin{ficha}{RF-011 — Shell único con carga de Frames sin recarga}
Descripción & La estructura general (barra lateral y superior) se carga una
sola vez; los Frames se muestran en un marco interno. \\ \hline
Prioridad & Alta \\ \hline
Actor & Usuario estándar, Administrador \\ \hline
Entradas & Selección de una opción del menú. \\ \hline
Proceso & El shell intercepta el clic y reemplaza la ubicación del marco
interno por la página del Frame, sin recargar la estructura ni acumular
historial. \\ \hline
Salidas & Frame visible dentro del shell. \\ \hline
Precondiciones & Sesión válida. \\ \hline
Postcondiciones & Barra lateral y superior permanecen intactas. \\ \hline
\end{ficha}

\begin{ficha}{RF-012 — Menú lateral según permisos y organización personal}
Descripción & La barra lateral muestra únicamente los ItemMenus activos y
permitidos al rol, organizados según los SuperMenus y el orden personal del
usuario. \\ \hline
Prioridad & Alta \\ \hline
Actor & Usuario estándar, Administrador \\ \hline
Entradas & Datos de menú provistos por el servidor. \\ \hline
Proceso & El servidor filtra por estado global y permiso de lectura del rol,
aplica la organización personal (SuperMenus propios, orden guardado) y
omite los SuperMenus vacíos; el cliente los renderiza como acordeones. \\ \hline
Salidas & Menú personalizado del usuario. \\ \hline
Precondiciones & Sesión válida. \\ \hline
Postcondiciones & --- \\ \hline
\end{ficha}

\begin{ficha}{RF-013 — Ítem activo y título dinámico}
Descripción & El shell resalta la opción de menú correspondiente al Frame
visible y muestra su título en la barra superior. \\ \hline
Prioridad & Media \\ \hline
Actor & Sistema \\ \hline
Entradas & Carga de un Frame en el marco. \\ \hline
Proceso & Al cargar el marco se lee el título del documento, se actualiza la
barra superior y la pestaña, y se marca el ítem cuyo destino coincide con la
ruta cargada (abriendo su acordeón contenedor). \\ \hline
Salidas & Ítem resaltado con el degradado corporativo; título actualizado. \\ \hline
Precondiciones & Frame cargado. \\ \hline
Postcondiciones & --- \\ \hline
\end{ficha}

\begin{ficha}{RF-014 — Enlace profundo por hash}
Descripción & La sección visible se refleja en el hash de la URL, de modo
que recargar la página conserva la sección actual. \\ \hline
Prioridad & Media \\ \hline
Actor & Sistema \\ \hline
Entradas & Navegación entre Frames; recarga del navegador. \\ \hline
Proceso & Cada carga del marco actualiza el hash; al iniciar, el shell carga
la ruta del hash si es una página interna válida del sistema, o el Dashboard
en caso contrario. \\ \hline
Salidas & Restauración de la sección tras recargar. \\ \hline
Precondiciones & Sesión válida. \\ \hline
Postcondiciones & --- \\ \hline
\end{ficha}

\begin{ficha}{RF-015 — Navegación móvil}
Descripción & En pantallas angostas la barra lateral se oculta y se despliega
con un botón hamburguesa sobre una capa oscura. \\ \hline
Prioridad & Media \\ \hline
Actor & Usuario estándar, Administrador \\ \hline
Entradas & Pulsación del botón hamburguesa o de la capa oscura. \\ \hline
Proceso & La barra se desliza desde el borde; seleccionar una opción o tocar
la capa la cierra. \\ \hline
Salidas & Menú accesible en dispositivos móviles. \\ \hline
Precondiciones & Ancho de pantalla igual o inferior a 768\,px. \\ \hline
Postcondiciones & --- \\ \hline
\end{ficha}

\begin{ficha}{RF-016 — Bloqueo de acceso directo no autorizado}
Descripción & Acceder por URL a un Frame sin permiso o con el módulo
deshabilitado muestra un aviso y redirige al Dashboard; los datos se
protegen además en el servidor. \\ \hline
Prioridad & Alta \\ \hline
Actor & Usuario estándar \\ \hline
Entradas & URL de un Frame no autorizado. \\ \hline
Proceso & El guardia del cliente comprueba el permiso de lectura y la
presencia del módulo en el menú vigente; en paralelo, todo endpoint valida
sesión, permiso y disponibilidad global antes de operar. \\ \hline
Salidas & Alerta ``No tiene permiso...'' y redirección; los endpoints
devuelven error. \\ \hline
Precondiciones & Sesión válida sin el permiso requerido. \\ \hline
Postcondiciones & Ningún dato del módulo es servido. \\ \hline
\end{ficha}

% ────────────────────────────────────────────────────────────
\section{Dashboard}

\begin{ficha}{RF-017 — Saludo contextual}
Descripción & El Dashboard saluda al usuario según la hora (buenos días /
buenas tardes / buenas noches) y muestra la fecha completa y su rol. \\ \hline
Prioridad & Baja \\ \hline
Actor & Usuario estándar, Administrador \\ \hline
Entradas & Hora local y datos de sesión. \\ \hline
Proceso & Se calcula el saludo por franja horaria y se formatea la fecha en
español. \\ \hline
Salidas & Encabezado personalizado con insignia del rol. \\ \hline
Precondiciones & Sesión válida. \\ \hline
Postcondiciones & --- \\ \hline
\end{ficha}

\begin{ficha}{RF-018 — Accesos rápidos según permisos}
Descripción & El Dashboard muestra botones de acceso directo a los Frames
permitidos al rol (excepto el propio Dashboard). \\ \hline
Prioridad & Media \\ \hline
Actor & Usuario estándar, Administrador \\ \hline
Entradas & Menú vigente del usuario. \\ \hline
Proceso & Se generan enlaces a partir de los ItemMenus con página,
descartando los ocultados por el usuario (RF-019). \\ \hline
Salidas & Botonera de accesos rápidos con ícono y nombre. \\ \hline
Precondiciones & Sesión válida. \\ \hline
Postcondiciones & --- \\ \hline
\end{ficha}

\begin{ficha}{RF-019 — Personalización de accesos rápidos}
Descripción & Cada usuario elige qué accesos rápidos visualizar, sin afectar
sus permisos. \\ \hline
Prioridad & Media \\ \hline
Actor & Usuario estándar, Administrador \\ \hline
Entradas & Selección de casillas en el modo ``Personalizar''. \\ \hline
Proceso & El botón Personalizar muestra todos los Frames permitidos con
casillas; al guardar, los desmarcados se registran como ocultos para ese
usuario (\texttt{accesos\_ocultos\_usuario}) y la botonera se actualiza. \\ \hline
Salidas & Confirmación y botonera filtrada. \\ \hline
Precondiciones & Sesión válida. \\ \hline
Postcondiciones & La preferencia persiste entre sesiones; el Frame sigue
accesible desde el menú. \\ \hline
\end{ficha}

\begin{ficha}{RF-020 — Estadísticas del sistema (solo Administrador)}
Descripción & El Dashboard del Administrador muestra el conteo exacto de
usuarios, roles y permisos registrados. \\ \hline
Prioridad & Media \\ \hline
Actor & Administrador \\ \hline
Entradas & --- \\ \hline
Proceso & El servidor cuenta las filas de las tablas correspondientes; el
cliente oculta la sección para cualquier rol distinto de Administrador. \\ \hline
Salidas & Tarjetas con los tres conteos. \\ \hline
Precondiciones & Sesión de Administrador. \\ \hline
Postcondiciones & --- \\ \hline
\end{ficha}

\begin{ficha}{RF-021 — Frase motivadora (solo usuarios estándar)}
Descripción & Bajo los accesos rápidos, los usuarios no administradores ven
una frase motivadora con su autor, seleccionada aleatoriamente de la base de
datos. \\ \hline
Prioridad & Baja \\ \hline
Actor & Usuario estándar \\ \hline
Entradas & --- \\ \hline
Proceso & El servidor devuelve una frase activa al azar; el cliente la
conserva durante la sesión (una frase por inicio de sesión) y la libera al
cerrar sesión. Si no hay frases, la sección no se muestra. \\ \hline
Salidas & Cita con formato ``frase'' — autor. \\ \hline
Precondiciones & Sesión de usuario no administrador; al menos una frase
activa. \\ \hline
Postcondiciones & --- \\ \hline
\end{ficha}

% ────────────────────────────────────────────────────────────
\section{Usuarios}

\begin{ficha}{RF-022 — Listar usuarios}
Descripción & Muestra el listado de cuentas con nombre, usuario, correo,
rol y estado. \\ \hline
Prioridad & Alta \\ \hline
Actor & Usuario estándar (solo lectura), Administrador \\ \hline
Entradas & --- \\ \hline
Proceso & El servidor devuelve los usuarios con su rol (unión con
\texttt{roles}) previa verificación del permiso de lectura. \\ \hline
Salidas & Tabla de usuarios con insignias de rol y estado. \\ \hline
Precondiciones & Permiso \texttt{usuarios:leer}. \\ \hline
Postcondiciones & --- \\ \hline
\end{ficha}

\begin{ficha}{RF-023 — Búsqueda de usuarios en vivo}
Descripción & Un campo de búsqueda filtra la tabla al escribir, sin recargar. \\ \hline
Prioridad & Media \\ \hline
Actor & Usuario estándar, Administrador \\ \hline
Entradas & Texto de búsqueda. \\ \hline
Proceso & Se comparan usuario, nombre, correo y rol de cada fila con el
texto, ocultando las no coincidentes. \\ \hline
Salidas & Tabla filtrada al instante. \\ \hline
Precondiciones & Listado cargado. \\ \hline
Postcondiciones & --- \\ \hline
\end{ficha}

\begin{ficha}{RF-024 — Filtros por rol y estado}
Descripción & El listado se filtra por rol (con opción ``Todos los roles'')
y por estado (Todos / Activos / Inactivos), de forma combinable entre sí y
con la búsqueda. \\ \hline
Prioridad & Media \\ \hline
Actor & Usuario estándar, Administrador \\ \hline
Entradas & Selección en los dos filtros. \\ \hline
Proceso & Las opciones de rol se construyen con los roles presentes en el
listado; al cambiar cualquier filtro la tabla se actualiza automáticamente. \\ \hline
Salidas & Tabla filtrada. \\ \hline
Precondiciones & Listado cargado. \\ \hline
Postcondiciones & --- \\ \hline
\end{ficha}

\begin{ficha}{RF-025 — Crear usuario}
Descripción & El Administrador registra cuentas con nombre, usuario, correo,
rol y contraseña inicial. \\ \hline
Prioridad & Alta \\ \hline
Actor & Administrador \\ \hline
Entradas & Nombre completo, nombre de usuario, correo, rol, contraseña. \\ \hline
Proceso & Se validan campos obligatorios, formato de correo, política de
contraseña y unicidad de usuario y correo; la contraseña se almacena con
bcrypt; la cuenta nace activa con la marca de primer inicio; se genera la
semilla del menú inicial (RF-043). \\ \hline
Salidas & Confirmación y retorno al listado. \\ \hline
Precondiciones & Permiso \texttt{usuarios:crear}; solo un Administrador
puede asignar el rol Admin. \\ \hline
Postcondiciones & Cuenta creada con menú inicial y cambio de contraseña
obligatorio pendiente. \\ \hline
\end{ficha}

\begin{ficha}{RF-026 — Editar usuario}
Descripción & Permite modificar nombre, usuario, correo, rol y,
opcionalmente, restablecer la contraseña. \\ \hline
Prioridad & Alta \\ \hline
Actor & Administrador \\ \hline
Entradas & Datos modificados; contraseña nueva opcional. \\ \hline
Proceso & Se validan unicidad y formato excluyendo al propio usuario; si se
proporciona contraseña, se aplica la política, se re-hashea y se reactiva la
marca de primer inicio. \\ \hline
Salidas & Confirmación y retorno al listado. \\ \hline
Precondiciones & Permiso \texttt{usuarios:editar}; no puede cambiarse el rol
del Administrador principal; solo un Administrador asigna el rol Admin. \\ \hline
Postcondiciones & Datos actualizados; si hubo contraseña nueva, el usuario
deberá cambiarla en su próximo acceso. \\ \hline
\end{ficha}

\begin{ficha}{RF-027 — Activar / desactivar usuario}
Descripción & Alterna el estado de una cuenta; una cuenta inactiva no puede
iniciar sesión. \\ \hline
Prioridad & Alta \\ \hline
Actor & Administrador \\ \hline
Entradas & Confirmación sobre el usuario elegido. \\ \hline
Proceso & Previa confirmación, el servidor invierte el estado; el
Administrador principal (usuario 1) no puede desactivarse. \\ \hline
Salidas & Mensaje ``Usuario activado/desactivado'' y listado actualizado. \\ \hline
Precondiciones & Permiso \texttt{usuarios:editar}. \\ \hline
Postcondiciones & El estado nuevo rige el próximo intento de login. \\ \hline
\end{ficha}

\begin{ficha}{RF-028 — Interfaz de usuarios condicionada por permisos}
Descripción & Los controles de escritura solo se muestran a quien posee la
acción correspondiente. \\ \hline
Prioridad & Alta \\ \hline
Actor & Usuario estándar, Administrador \\ \hline
Entradas & Permisos del rol en sesión. \\ \hline
Proceso & El botón ``Nuevo Usuario'' se oculta sin \texttt{usuarios:crear};
los botones Editar y Activar/Desactivar solo aparecen con
\texttt{usuarios:editar}; el servidor revalida cada operación. \\ \hline
Salidas & Vista de solo lectura para usuarios estándar. \\ \hline
Precondiciones & Listado cargado. \\ \hline
Postcondiciones & --- \\ \hline
\end{ficha}

% ────────────────────────────────────────────────────────────
\section{Roles}

\begin{ficha}{RF-029 — Listar roles}
Descripción & Muestra los roles con nombre, descripción y estado, con
búsqueda en vivo. \\ \hline
Prioridad & Alta \\ \hline
Actor & Administrador \\ \hline
Entradas & Texto de búsqueda opcional. \\ \hline
Proceso & El servidor devuelve los roles previa verificación del permiso;
el rol Admin no ofrece botón de edición. \\ \hline
Salidas & Tabla de roles. \\ \hline
Precondiciones & Permiso \texttt{roles:leer} (solo lo posee el
Administrador). \\ \hline
Postcondiciones & --- \\ \hline
\end{ficha}

\begin{ficha}{RF-030 — Crear rol}
Descripción & Registra un rol nuevo con nombre único, descripción opcional y
estado. \\ \hline
Prioridad & Media \\ \hline
Actor & Administrador \\ \hline
Entradas & Nombre del rol, descripción, estado. \\ \hline
Proceso & Se valida el nombre obligatorio y su unicidad; el rol se crea sin
permisos (deben asignarse en el módulo Permisos). \\ \hline
Salidas & Confirmación y retorno al listado. \\ \hline
Precondiciones & Permiso \texttt{roles:crear}. \\ \hline
Postcondiciones & Rol disponible para asignar a usuarios y configurar
permisos. \\ \hline
\end{ficha}

\begin{ficha}{RF-031 — Editar rol}
Descripción & Modifica nombre, descripción y estado de un rol existente. \\ \hline
Prioridad & Media \\ \hline
Actor & Administrador \\ \hline
Entradas & Datos modificados del rol. \\ \hline
Proceso & Se valida la unicidad del nombre excluyendo al propio rol;
desactivar un rol impide el acceso de todos sus usuarios (RF-010). \\ \hline
Salidas & Confirmación y retorno al listado. \\ \hline
Precondiciones & Permiso \texttt{roles:editar}; el rol Admin no es editable
desde la interfaz. \\ \hline
Postcondiciones & Cambios aplicados a los próximos inicios de sesión. \\ \hline
\end{ficha}

% ────────────────────────────────────────────────────────────
\section{Permisos}

\begin{ficha}{RF-032 — Selección del rol a configurar}
Descripción & El módulo Permisos presenta un selector con todos los roles
excepto Admin, cuyo acceso total es implícito y no configurable. \\ \hline
Prioridad & Alta \\ \hline
Actor & Administrador \\ \hline
Entradas & Rol elegido. \\ \hline
Proceso & Al seleccionar un rol se cargan sus permisos vigentes y se
construyen las dos matrices (RF-033 y RF-034). \\ \hline
Salidas & Matrices de permisos del rol. \\ \hline
Precondiciones & Permiso \texttt{permisos:leer}. \\ \hline
Postcondiciones & --- \\ \hline
\end{ficha}

\begin{ficha}{RF-033 — Matriz de acceso a Frames}
Descripción & Primera matriz: todos los Frames del sistema con una casilla
``Habilitar'' que otorga o revoca el acceso (acción \texttt{leer}). \\ \hline
Prioridad & Alta \\ \hline
Actor & Administrador \\ \hline
Entradas & Casillas marcadas o desmarcadas. \\ \hline
Proceso & Los Frames provienen del registro central del servidor; habilitar
un Frame lo hace visible en el menú del rol y accesible; deshabilitarlo lo
oculta y bloquea por completo. \\ \hline
Salidas & Estado de acceso por Frame. \\ \hline
Precondiciones & Rol seleccionado. \\ \hline
Postcondiciones & Cambios pendientes hasta guardar (RF-036). \\ \hline
\end{ficha}

\begin{ficha}{RF-034 — Acciones internas por Frame habilitado}
Descripción & Segunda sección, generada dinámicamente: por cada Frame
habilitado que declare acciones internas se listan sus casillas
(p.~ej. Usuarios: agregar, editar/desactivar; Frame~1: guardar, editar,
eliminar tarea). \\ \hline
Prioridad & Alta \\ \hline
Actor & Administrador \\ \hline
Entradas & Casillas de acciones internas. \\ \hline
Proceso & Al alternar la casilla de un Frame, su bloque de acciones aparece
o desaparece en el acto; las acciones de Frames deshabilitados no se
muestran ni se evalúan al guardar. \\ \hline
Salidas & Configuración fina de operaciones por rol. \\ \hline
Precondiciones & Frame habilitado con acciones declaradas en el registro. \\ \hline
Postcondiciones & Cambios pendientes hasta guardar (RF-036). \\ \hline
\end{ficha}

\begin{ficha}{RF-035 — Marcar / desmarcar todos}
Descripción & Botones que activan o desactivan de una vez todos los Frames y
todas sus acciones internas. \\ \hline
Prioridad & Baja \\ \hline
Actor & Administrador \\ \hline
Entradas & Pulsación del botón correspondiente. \\ \hline
Proceso & Se actualizan todas las casillas y la visibilidad de los bloques
de acciones. \\ \hline
Salidas & Matrices completamente marcadas o vacías. \\ \hline
Precondiciones & Rol seleccionado. \\ \hline
Postcondiciones & Cambios pendientes hasta guardar. \\ \hline
\end{ficha}

\begin{ficha}{RF-036 — Guardar permisos}
Descripción & Persiste la configuración completa del rol seleccionado. \\ \hline
Prioridad & Alta \\ \hline
Actor & Administrador \\ \hline
Entradas & Estado de ambas matrices. \\ \hline
Proceso & El cliente envía los pares módulo--acción de los Frames
habilitados; el servidor valida cada par contra el registro central,
reemplaza los permisos del rol dentro de una transacción y rechaza
modificaciones al rol Admin. \\ \hline
Salidas & Confirmación ``Permisos actualizados correctamente''. \\ \hline
Precondiciones & Permiso \texttt{permisos:crear}; rol distinto de Admin. \\ \hline
Postcondiciones & Los usuarios del rol reciben los permisos nuevos en su
próxima verificación de sesión. \\ \hline
\end{ficha}

% ────────────────────────────────────────────────────────────
\section{Menú (personalización individual)}

\begin{ficha}{RF-037 — Organizador del menú personal}
Descripción & Todo usuario dispone de un organizador visual de su menú:
un árbol libre donde ItemMenus y SuperMenus propios se intercalan en la
raíz. \\ \hline
Prioridad & Alta \\ \hline
Actor & Usuario estándar, Administrador \\ \hline
Entradas & --- \\ \hline
Proceso & El servidor entrega los ItemMenus activos permitidos al rol (sin
URLs) junto con los SuperMenus y la organización del usuario; el organizador
los presenta como lista arrastrable. \\ \hline
Salidas & Vista editable del menú personal. \\ \hline
Precondiciones & Permiso \texttt{menu:leer}. \\ \hline
Postcondiciones & --- \\ \hline
\end{ficha}

\begin{ficha}{RF-038 — Crear SuperMenu}
Descripción & El usuario crea contenedores propios para agrupar sus
ItemMenus. \\ \hline
Prioridad & Alta \\ \hline
Actor & Usuario estándar, Administrador \\ \hline
Entradas & Nombre del SuperMenu. \\ \hline
Proceso & Se solicita el nombre; el servidor lo registra a nombre del
usuario al final de su menú. \\ \hline
Salidas & SuperMenu vacío listo para recibir ítems. \\ \hline
Precondiciones & Permiso \texttt{menu:leer}; nombre no vacío. \\ \hline
Postcondiciones & El SuperMenu pertenece exclusivamente a ese usuario. \\ \hline
\end{ficha}

\begin{ficha}{RF-039 — Renombrar SuperMenu}
Descripción & El usuario cambia el nombre de un SuperMenu propio. \\ \hline
Prioridad & Media \\ \hline
Actor & Usuario estándar, Administrador \\ \hline
Entradas & Nuevo nombre. \\ \hline
Proceso & El servidor actualiza el nombre validando la propiedad del
SuperMenu. \\ \hline
Salidas & Confirmación y organizador actualizado. \\ \hline
Precondiciones & SuperMenu propio; nombre no vacío. \\ \hline
Postcondiciones & --- \\ \hline
\end{ficha}

\begin{ficha}{RF-040 — Eliminar SuperMenu vacío}
Descripción & Solo pueden eliminarse SuperMenus sin ItemMenus dentro. \\ \hline
Prioridad & Media \\ \hline
Actor & Usuario estándar, Administrador \\ \hline
Entradas & Confirmación de eliminación. \\ \hline
Proceso & El botón Eliminar solo se muestra en SuperMenus vacíos; el
servidor verifica propiedad y vacuidad antes de borrar. \\ \hline
Salidas & Confirmación, o rechazo si contiene elementos. \\ \hline
Precondiciones & SuperMenu propio y vacío. \\ \hline
Postcondiciones & Los ItemMenus del sistema no se ven afectados. \\ \hline
\end{ficha}

\begin{ficha}{RF-041 — Reorganización por arrastrar y soltar}
Descripción & El usuario reordena ItemMenus y SuperMenus, mueve ítems entre
SuperMenus y los intercala libremente en la raíz. El Dashboard puede
reordenarse pero nunca agruparse. \\ \hline
Prioridad & Alta \\ \hline
Actor & Usuario estándar, Administrador \\ \hline
Entradas & Gestos de arrastrar y soltar. \\ \hline
Proceso & Soltar un elemento sobre otro nodo de la raíz lo inserta en esa
posición; soltar un ítem dentro de un SuperMenu lo agrupa; intentar agrupar
el Dashboard muestra un aviso y lo devuelve a la raíz. \\ \hline
Salidas & Árbol reorganizado en pantalla. \\ \hline
Precondiciones & Organizador cargado. \\ \hline
Postcondiciones & Cambio persistido automáticamente (RF-042). \\ \hline
\end{ficha}

\begin{ficha}{RF-042 — Persistencia y refresco del menú personal}
Descripción & Cada movimiento se guarda automáticamente; el botón Guardar
persiste, confirma y refresca el menú lateral del shell sin cerrar sesión. \\ \hline
Prioridad & Alta \\ \hline
Actor & Usuario estándar, Administrador \\ \hline
Entradas & Organización vigente del árbol. \\ \hline
Proceso & Se envía la secuencia completa (orden de raíz compartido entre
ítems y SuperMenus, orden interno por SuperMenu); el servidor la registra en
una transacción validando la propiedad; el botón Guardar recarga además la
barra lateral del shell. \\ \hline
Salidas & ``Menú guardado y actualizado''; sidebar reflejando el cambio. \\ \hline
Precondiciones & Permiso \texttt{menu:leer}. \\ \hline
Postcondiciones & La organización se reconstruye idéntica en próximos
inicios de sesión; la de otros usuarios no cambia. \\ \hline
\end{ficha}

\begin{ficha}{RF-043 — Menú inicial automático (semilla)}
Descripción & Al crear una cuenta se genera una única vez su estructura de
menú: Dashboard primero, SuperMenu ``Mi Perfil'' (Ver Perfil, Cambiar
Contraseña, Menú) y el resto sin agrupar; para administradores se añaden los
SuperMenus ``Administración'' y ``Frames''. \\ \hline
Prioridad & Alta \\ \hline
Actor & Sistema \\ \hline
Entradas & Creación de un usuario (RF-025). \\ \hline
Proceso & La semilla consulta los ItemMenus activos permitidos al rol, crea
los SuperMenus predefinidos con su contenido en el orden establecido y ubica
el resto en la raíz; se aborta si el usuario ya posee organización. \\ \hline
Salidas & Menú inicial completo en el primer acceso. \\ \hline
Precondiciones & Usuario recién creado sin organización previa. \\ \hline
Postcondiciones & El sistema no vuelve a reconstruir el menú jamás; toda
personalización posterior se conserva. \\ \hline
\end{ficha}

% ────────────────────────────────────────────────────────────
\section{Configurar Menús (disponibilidad global)}

\begin{ficha}{RF-044 — Listado global de ItemMenus}
Descripción & El Administrador visualiza todos los ItemMenus del sistema con
su módulo y disponibilidad, ordenados por identificador ascendente. \\ \hline
Prioridad & Alta \\ \hline
Actor & Administrador \\ \hline
Entradas & --- \\ \hline
Proceso & El servidor lista la tabla \texttt{menu} completa previa
verificación del permiso \texttt{configmenu:leer}. \\ \hline
Salidas & Tabla estable con casilla ``Activo'' por ItemMenu. \\ \hline
Precondiciones & Sesión de Administrador. \\ \hline
Postcondiciones & --- \\ \hline
\end{ficha}

\begin{ficha}{RF-045 — Activar / desactivar ItemMenus globalmente}
Descripción & Las casillas acumulan cambios que el botón Guardar persiste en
lote, con confirmación y refresco inmediato del menú lateral. \\ \hline
Prioridad & Alta \\ \hline
Actor & Administrador \\ \hline
Entradas & Casillas modificadas. \\ \hline
Proceso & Al guardar se envía cada cambio al servidor; ``Configurar Menús''
no puede desactivarse (su casilla está deshabilitada y el servidor lo
rechaza); al terminar se recarga la tabla y el sidebar del shell. \\ \hline
Salidas & ``Cambios guardados (N ItemMenu actualizados)''. \\ \hline
Precondiciones & Permiso \texttt{configmenu:editar}. \\ \hline
Postcondiciones & La disponibilidad rige de inmediato para todos los
usuarios. \\ \hline
\end{ficha}

\begin{ficha}{RF-046 — Bloqueo integral de módulos deshabilitados}
Descripción & Un ItemMenu inactivo desaparece del menú de todos los usuarios
y su módulo queda inaccesible por navegación, URL directa y endpoints. \\ \hline
Prioridad & Alta \\ \hline
Actor & Sistema \\ \hline
Entradas & Estado global de los ItemMenus. \\ \hline
Proceso & El listado del menú filtra por estado; el guardia del cliente
bloquea módulos ausentes del menú vigente; cada endpoint verifica la
disponibilidad global del módulo y responde ``deshabilitado temporalmente''
incluso al Administrador. \\ \hline
Salidas & Módulo completamente inaccesible mientras esté inactivo. \\ \hline
Precondiciones & ItemMenu desactivado en Configurar Menús. \\ \hline
Postcondiciones & Reactivarlo restaura el acceso sin pérdida de permisos ni
de organización personal. \\ \hline
\end{ficha}

% ────────────────────────────────────────────────────────────
\section{Mi Perfil}

\begin{ficha}{RF-047 — Ver perfil propio}
Descripción & Muestra exclusivamente los datos del usuario autenticado:
nombre, usuario, correo, rol y fecha de registro. \\ \hline
Prioridad & Media \\ \hline
Actor & Usuario estándar, Administrador \\ \hline
Entradas & --- \\ \hline
Proceso & El servidor consulta la fila del usuario en sesión (unida a su
rol); ningún parámetro permite consultar a terceros. \\ \hline
Salidas & Ficha de datos personales. \\ \hline
Precondiciones & Permiso \texttt{perfil:leer}. \\ \hline
Postcondiciones & --- \\ \hline
\end{ficha}

\begin{ficha}{RF-048 — Cambiar contraseña}
Descripción & El usuario cambia su contraseña validando la actual y
cumpliendo la política de seguridad, con indicador de fortaleza y botones de
ver/ocultar por campo. \\ \hline
Prioridad & Alta \\ \hline
Actor & Usuario estándar, Administrador \\ \hline
Entradas & Contraseña actual, nueva y su confirmación. \\ \hline
Proceso & Cliente y servidor validan: actual correcta, política (mínimo 8
caracteres, una mayúscula, un número, un carácter especial), confirmación
coincidente y nueva distinta de la actual; se almacena el hash bcrypt y se
desactiva la marca de primer inicio. \\ \hline
Salidas & Confirmación y redirección al Dashboard. \\ \hline
Precondiciones & Sesión válida. \\ \hline
Postcondiciones & La contraseña anterior deja de ser válida. \\ \hline
\end{ficha}

% ────────────────────────────────────────────────────────────
\section{Frame 1 — Tareas (CRUD de ejemplo)}

\begin{ficha}{RF-049 — Listar tareas}
Descripción & Muestra las tareas registradas con título, descripción y
nombre del creador. \\ \hline
Prioridad & Media \\ \hline
Actor & Usuario estándar, Administrador \\ \hline
Entradas & --- \\ \hline
Proceso & El servidor devuelve las tareas (unidas al nombre del creador)
previa verificación de \texttt{frame1:leer}. \\ \hline
Salidas & Tabla de tareas. \\ \hline
Precondiciones & Permiso \texttt{frame1:leer}. \\ \hline
Postcondiciones & --- \\ \hline
\end{ficha}

\begin{ficha}{RF-050 — Crear tarea}
Descripción & Registra una tarea con título obligatorio y descripción
opcional, asociada a su creador. \\ \hline
Prioridad & Media \\ \hline
Actor & Usuario estándar, Administrador \\ \hline
Entradas & Título, descripción. \\ \hline
Proceso & Validación en cliente y servidor del permiso \texttt{frame1:crear}
y del título; se registra el usuario creador. \\ \hline
Salidas & Confirmación y listado actualizado. \\ \hline
Precondiciones & Permiso \texttt{frame1:crear}. \\ \hline
Postcondiciones & Tarea visible para todos los usuarios con acceso. \\ \hline
\end{ficha}

\begin{ficha}{RF-051 — Editar tarea}
Descripción & Modifica el título y la descripción de una tarea existente. \\ \hline
Prioridad & Media \\ \hline
Actor & Usuario estándar, Administrador \\ \hline
Entradas & Tarea seleccionada; datos modificados. \\ \hline
Proceso & El botón Editar precarga el formulario en modo actualización, con
opción de cancelar; el servidor verifica \texttt{frame1:editar}. \\ \hline
Salidas & Confirmación y listado actualizado. \\ \hline
Precondiciones & Permiso \texttt{frame1:editar}. \\ \hline
Postcondiciones & --- \\ \hline
\end{ficha}

\begin{ficha}{RF-052 — Eliminar tarea}
Descripción & Elimina definitivamente una tarea, previa confirmación. \\ \hline
Prioridad & Media \\ \hline
Actor & Usuario estándar (si su rol lo permite), Administrador \\ \hline
Entradas & Confirmación sobre la tarea elegida. \\ \hline
Proceso & El servidor verifica \texttt{frame1:eliminar} y borra la fila. \\ \hline
Salidas & Confirmación y listado actualizado. \\ \hline
Precondiciones & Permiso \texttt{frame1:eliminar}. \\ \hline
Postcondiciones & La tarea no puede recuperarse. \\ \hline
\end{ficha}

\begin{ficha}{RF-053 — Interfaz de tareas condicionada por permisos}
Descripción & El Frame~1 adapta su interfaz a las acciones del rol: sin
crear ni editar se oculta el formulario; los botones Editar y Eliminar solo
aparecen con su permiso. \\ \hline
Prioridad & Alta \\ \hline
Actor & Usuario estándar \\ \hline
Entradas & Permisos del rol en sesión. \\ \hline
Proceso & El cliente consulta los permisos antes de renderizar; el servidor
revalida cada operación, de modo que la restricción no depende de la
interfaz. \\ \hline
Salidas & Ejemplo funcional de permisos por operación (p.~ej. Rol1 crea y
edita pero no elimina; Rol2 dispone del CRUD completo; Rol3 no accede). \\ \hline
Precondiciones & Frame~1 habilitado para el rol. \\ \hline
Postcondiciones & --- \\ \hline
\end{ficha}

\begin{ficha}{RF-054 — Frames de demostración 2 a 5}
Descripción & Cuatro pantallas de ejemplo que muestran el texto ``Aquí va un
Frame'' y demuestran la distribución de accesos por rol. \\ \hline
Prioridad & Baja \\ \hline
Actor & Usuario estándar, Administrador \\ \hline
Entradas & Navegación al Frame. \\ \hline
Proceso & Cada Frame verifica sesión y permiso de lectura de su módulo
(\texttt{frame2}\dots\texttt{frame5}) antes de mostrarse; la distribución
vigente asigna a cada rol al menos un Frame y ningún Frame a todos los
roles. \\ \hline
Salidas & Contenido de demostración. \\ \hline
Precondiciones & Permiso \texttt{frameN:leer}. \\ \hline
Postcondiciones & --- \\ \hline
\end{ficha}

% ============================================================
% CAPÍTULO 4 — REQUERIMIENTOS NO FUNCIONALES
% Fichas RNF-001 ... RNF-0NN por categoría:
%   Seguridad, Rendimiento, Usabilidad, Disponibilidad,
%   Confiabilidad, Persistencia, Integridad, Compatibilidad,
%   Arquitectura, Mantenibilidad
% ============================================================
\chapter{Requerimientos no funcionales}\label{cap:rnf}

\textit{El contenido de este capítulo se desarrolla en la Fase 3 de la documentación.}

% ============================================================
% CAPÍTULO 5 — CASOS DE USO
% Fichas CU-01 ... CU-NN con: Nombre, Actor, Descripción,
% Precondiciones, Flujo principal, Flujos alternativos,
% Postcondiciones. Agrupados por actor/módulo.
% ============================================================
\chapter{Casos de uso}\label{cap:cu}

\textit{El contenido de este capítulo se desarrolla en la Fase 3 de la documentación.}

% ============================================================
% CAPÍTULO 6 — REGLAS DE NEGOCIO
% Fichas RN-01 ... RN-23 por categoría:
%   Autenticación y sesión / RBAC / Menús / Dashboard
% Cada regla con su evidencia (archivo del código que la implementa).
% ============================================================
\chapter{Reglas de negocio}\label{cap:rn}

\textit{El contenido de este capítulo se desarrolla en la Fase 3 de la documentación.}

% ============================================================
% CAPÍTULO 7 — ANEXOS
% Previsto:
%   A. Modelo de datos (11 tablas, columnas y relaciones)
%   B. Matriz de permisos por rol (estado vigente)
%   C. Catálogo de endpoints del servidor
%   D. Estructura de archivos del proyecto
% ============================================================
\chapter{Anexos}\label{cap:anexos}

\textit{El contenido de este capítulo se desarrolla en la Fase 3 de la documentación.}


\end{document}
